\subsection{E1: literal counterexample and the relative-normal result}

Over $R=\mathbf C[[t]]$ take the plane cubic
\[
 X:\ Y^2Z=X^3+X^2Z+t^2Z^3,
 \quad O=[0:1:0],\quad P=[0:t:1].
\]
Its generic fibre is smooth since the cubic discriminant is
$-t^2(4+27t^2)$. Its special fibre is an irreducible nodal cubic. Its total
space is normal, because it is a hypersurface with an isolated singularity
at $(x,y,t)=(0,0,0)$. It is proper and flat, has a section, and has
$f_*\mathcal O_X=\mathcal O_R$ universally by the plane-cubic exact sequence.
The generic line bundle $\mathcal L_\eta=\mathcal O(P_\eta-O_\eta)$ is
algebraically equivalent to zero.

The closure $P$ has local ideal $(x,y-t)$ at the node, with two independent
linear generators, so it is not Cartier. If $\mathcal L_\eta$ extended to
an invertible sheaf on $X$, choose a rational trivialization agreeing with
the generic divisor $P_\eta-O_\eta$. Its divisor differs from $P-O$ by a
vertical divisor. The only vertical prime is the reduced irreducible special
fibre, whose class is the principal divisor of $t$. As $O$ is Cartier, this
would make $P$ Cartier, a contradiction. Thus mere total-space normality does
not suffice, even with projectivity.

For the original relative-normal reading, assume instead that $f$ is proper
flat with geometrically normal, geometrically connected fibres over a trait,
and has a section. Universal $H^0=\mathcal O$ gives the algebraic Picard space
by Artin, and geometrically integral fibres give separatedness. Raynaud
\cite[Theorem~3.3.1, p.~40]{Raynaud1970} makes this separated group algebraic
space a scheme. Grothendieck \cite[Theorem~1.1(i), Corollary~2.3]{Grothendieck1962VI}
then says that the union of its fibrewise neutral components is a proper
\emph{subset}, in his precise topological/valuative sense. Thus a generic
algebraically trivial Picard point extends uniquely to a Picard point with
special fibre algebraically trivial; the section supplies an actual line
bundle representing it. No projectivity is required.

It is essential not to recast this as representability and properness of
the fibrewise $\operatorname{Pic}^0$ subfunctor. That subfunctor need not be an
algebraic space, even for smooth projective families, as
\cite[Proposition~2.2]{LaurentSchroer2024} demonstrates for Enriques surfaces.
The topological subset statement and the representability statement differ.

\subsection{E2/H9: separatedness of correspondences}

Let $S$ be locally noetherian and $X,Y$ proper flat and finitely presented
over $S$, both with universal $H^0=\mathcal O_S$. Define $\operatorname{Corr}_S(X,Y)$
as the fppf quotient of $\operatorname{Pic}_{X\times_SY/S}$ by the two factors'
Picard functors. It is a separated unramified $S$-scheme locally of finite type.

Here are the valuative argument and the exact representability deduction.
Given a trait $T=\operatorname{Spec}R$ over $S$ and a correspondence class
vanishing over its generic point, make a faithfully flat extension of traits
so that $X_T,Y_T$ have sections $x,y$. Such sections are obtained by taking
closed points of the generic fibres over a finite field extension and using
properness over a dominating valuation ring. A correspondence is now uniquely
represented by a line bundle bi-rigidified along $X\times y$ and $x\times Y$:
from a representative $L$ one replaces it by
\[
 L\otimes p_X^*(L|_{X\times y})^{-1}
  \otimes p_Y^*(L|_{x\times Y})^{-1}
  \otimes f^*(L|_{(x,y)}).
\]
This bundle is trivial on the generic fibre. It defines a map
$a:X_T\to P=\operatorname{Pic}_{Y_T/T}$ with $a(x)=0$ and $a_\eta=0$.

Let $E\subset P$ be the schematic closure of the generic identity.
Raynaud \cite[\S3.2(c), p.~40]{Raynaud1970} defines this subsheaf by precisely
the maps from flat $T$-schemes whose generic maps land in the identity.
Since $X_T$ is flat, $a$ factors through $E$. Universal $H^0=\mathcal O_T$
for $Y_T$ and Raynaud's Proposition~5.2, p.~46, imply that $E/T$ is an
\emph{\'etale algebraic space}. Each geometric fibre of $X_T$ is connected,
and a map from it to the discrete fibre of $E$ is constant. The section fixes
that value as zero. The equality locus of $a$ and zero is open, since $E/T$
is unramified, and contains every geometric point; it is therefore all of
$X_T$. The bi-rigidified line bundle is trivial. Faithfully flat descent
returns the original correspondence class to zero.

This proves the separatedness condition (F6) in Murre
\cite[p.~245]{Murre1966}. Murre's Theorem~4, p.~256, proves that, under exactly
these hypotheses, \emph{if} the correspondence functor is separated then it
is represented by an unramified scheme locally of finite type. The conditional
theorem becomes applicable after the valuative proof just supplied; it must
not be cited as if it already included unconditional separatedness.

\subsection{E3 and H6: flattening and the numerical quotient}

For $f:X\to S$ proper of finite presentation and $F$ finitely presented and
quasi-coherent on $X$, the functor taking $T/S$ to a singleton exactly when
$F_T$ is flat over $T$ is represented by a monomorphism $S_F\to S$ of finite
presentation. This is precisely the requested Problem A: the scheme version
is Stacks, Lemma~38.27.6(3), tag~05PS; the algebraic-space version is
Lemma~77.11.8(3), tag~0CWQ. No projectivity hypothesis is used.

Suppose $P=\operatorname{Pic}_{X/S}$ is an algebraic space locally of finite
presentation and $H=P^\tau$ is flat, of finite presentation, open and closed.
The free translation action gives an algebraic-space quotient $Q=P/H$ with
$P\to Q$ an fppf torsor (Stacks, Lemma~80.11.7, tag~06PH). Its relation
$P\times H\to P\times P$ is open and closed. Faithfully flat descent of the
diagonal makes $Q/S$ unramified and separated. A separated locally
quasi-finite algebraic space is a scheme (Stacks, Lemma~67.51.1, tag~0418).
This proves both readings of H6. The prompt summary asks for an algebraic-space
quotient under flatness: that already follows from tag~06PH, without closedness.
The fuller wording in \texttt{HUNCHES.md} assumes closedness as well, and obtains
the stronger separated scheme conclusion proved above. Compare
\cite[Theorem~3.1]{LaurentSchroer2024}.

\subsection{Some decisive leaf-by-leaf answers on 57:308--315}

The accompanying CSV retains every substantive leaf, including duplicates and
ambiguous readings; 57:308 is a heading and 57:312 is an unrelated tensor-product
insertion. The remaining broad entries are not silently marked settled.

For $X:\ xy=tz^2\subset\mathbf P^2_{\mathbf C[[t]]}$, its relative Picard
scheme is the disjoint union over total degrees $n$ of schemes $S_n$ obtained
by gluing copies of $\operatorname{Spec}\mathbf C[[t]]$, indexed by pairs
$(a,b)$ with $a+b=n$, along the common generic point. This is the explicit
split version of \cite[Example~4.14]{Kleiman2005}. Thus a connected component
need not be of finite type; the closure of a single generic point need not
be quasi-compact; and $P^\tau$ (the identity copy inside $S_0$) need not be
closed. These answer 57:309(a),(a'),(a'') in general and 57:310(8) negatively.

For the cuspidal cubic $C:y^2z=x^3$ in characteristic $p>0$, normalization
and the conductor exact sequence give
$\operatorname{Pic}^0_C=\operatorname{Pic}^\tau_C=\mathbf G_a$:
the relevant unit quotient is
$(k[\epsilon]/\epsilon^2)^\times/k^\times\simeq\mathbf G_a$.
Multiplication by $p$ is the constant zero map and is not open. Every
$k$-point is killed by $p$, so the set of points killed by a fixed integer
need not be finite. This answers 57:310(4),(6) in the generality actually
written there. For an integer invertible on the base, multiplication on a
commutative group scheme locally of finite type is \`etale
\cite[Theorem~2.5]{Grothendieck1962VI}; in particular it is universally open.

The tau-part is always an open algebraic subspace of finite presentation
under the proper-flat-finitely-presented, cohomologically-flat-in-degree-zero
hypotheses: \cite[Theorem~2.1]{LaurentSchroer2024}. Its finite-type and bounded
NS-rank assertions have the classical sources SGA~6 XIII, Theorems~4.7 and
5.1; see the exact discussion in \cite[Remark~6.19]{Kleiman2005}. Multiplication
by every nonzero integer on a represented relative Picard functor is of finite
type by SGA~6 XIII, Theorem~3.6; \cite[Remark~6.28]{Kleiman2005} records the
full proper version. These finiteness theorems do not imply flatness or
universal openness of the tau-part.

For geometrically integral locally factorial fibres the tau-part is proper.
This follows from \cite[Proposition~2.3(ii)]{LaurentSchroer2024}. Under mere geometrically
normal fibres, this particular theorem has an extra locally factorial
hypothesis and cannot be cited after deleting it. The properness of the
\emph{neutral subset} in E1 similarly does not itself prove properness of
the entire tau-part.

\subsection{E4: separation under reduced fibres, and a represented counterexample}

The two questions on 57:313 are:
\begin{quote}
``S\'epar\'e sur $S$ si $X_s$ sont g\'eom$^t$ loc. int\`egres ?''
and ``$\underline{\mathrm{Pic}}^\tau_{X/S}$ s\'epar\'e sur $S$ si les $X_s$
s\'eparables ?''
\end{quote}
Both have affirmative answers in the proper-flat setting. For the second,
test separatedness after base change to a strictly henselian trait. Geometric
reducedness of the special fibre implies Raynaud's property (N)
\cite[Example~6.1.6, p.~49]{Raynaud1970}. After the Stein-factor reduction
of \S6.1.9 one has property (N)$^*$, and every special-fibre multiplicity is
one. Thus $d=1$ and $d'=1$ in Proposition~6.1.11. Theorem~6.4.1(3), p.~56,
or equivalently Corollary~6.4.5, p.~57, makes $P^\tau$ separated. This holds
after every trait base change, proving the valuative criterion over $S$.
Every $n$-torsion Picard point belongs to $P^\tau$, so the $n$-torsion
subfunctor is separated as well. For the first question, a proper geometrically
locally integral fibre is the disjoint union of integral connected components.
The finite \'etale Stein-factor reduction therefore reduces to the usual
separatedness theorem for geometrically integral fibres. Separatedness
descends along that finite \'etale cover.

Representability alone does not make $P^\tau$ separated, answering the more
general question on 57:309(e) negatively. Here is a complete example over
$\mathbf C$. Put $\lambda_i=i$ for $0\leq i\leq5$ and define
\[
 Y=\left\{\sum_{i=0}^5\lambda_i^jX_i^2=0\quad(j=0,1,2)\right\}
 \subset\mathbf P^5_{\mathbf C}.
\]
This is a smooth connected complete intersection of three quadrics. To check
smoothness, a nonzero dependence among the three Jacobian rows gives a nonzero
polynomial of degree at most two vanishing at every $\lambda_i$ for which
$X_i\ne0$. There would be at most two such coordinates; the corresponding
Vandermonde columns are independent, and the three defining equations force
all their squares to vanish, a contradiction. The Koszul complex computes
$H^0(Y,\mathcal O_Y)=\mathbf C$ and $H^1(Y,\mathcal O_Y)=0$; adjunction gives
$\omega_Y\simeq\mathcal O_Y$. The involution
\[
 g[X_0:X_1:X_2:X_3:X_4:X_5]
 =[X_0:X_1:X_2:-X_3:-X_4:-X_5]
\]
is free: its fixed locus consists of the two coordinate planes, and on each
the relevant three Vandermonde columns are independent, so it misses $Y$.
All twenty three-column determinants, including the two fixed-plane ones,
are computed exactly in the accompanying code.

Set $R=\mathbf C[[t]]$, $R'=\mathbf C[[u]]$ with $t=u^2$, and form the
projective quotient
\[
 \mathcal X=(Y\times\operatorname{Spec}R')/
 \langle(g,u\mapsto-u)\rangle\longrightarrow\operatorname{Spec}R.
\]
The action is free, so the quotient map is finite \'etale. Consequently
$\mathcal X$ is regular and flat over $R$. Write $Z=Y/\langle g\rangle$.
The special fibre is $2Z$. If $T$ is the nontrivial order-two line bundle
given by the cover, then
$\nu_*\mathcal O_Y=\mathcal O_Z\oplus T$ and the square-zero ideal in
$\mathcal X_s$ is $T$. The cover is connected, so $T$ is nontrivial and
$H^0(Z,T)=0$. The vanishing of $H^1(Y,\mathcal O_Y)$ gives
$H^1(Z,\mathcal O_Z)=H^1(Z,T)=0$. Hence
$H^0(\mathcal X_s,\mathcal O)=\mathbf C$ and $H^1(\mathcal X_s,\mathcal O)=0$.
The same statements hold geometrically on the generic fibre, a form of $Y$.
Thus $f_*\mathcal O=\mathcal O_R$ universally and the Picard algebraic space
is unramified. Its identity is an open subgroup represented by $\operatorname{Spec}R$;
Raynaud's Proposition~5.3, since the residue field is separably closed, makes
the entire Picard space a scheme.

Nevertheless $\mathcal L=\mathcal O_{\mathcal X}(Z)$ satisfies
$\mathcal L^{\otimes2}\simeq\mathcal O_{\mathcal X}$ and
$\mathcal L_\eta\simeq\mathcal O$, while $\mathcal L|_Z\simeq T\ne\mathcal O_Z$.
Here $Z$ is Cartier because $\mathcal X$ is regular, and $2Z=\operatorname{div}(t)$.
So $\mathcal L$ and the identity are two sections of the represented tau-part
that agree generically and disagree on the special fibre. The tau-part is not
separated. The special fibre in this example is nonreduced; it does not
contradict the reduced-fibre theorem.

\subsection{E4: universal openness and constant dimension fail in general}

The repeated question ``$\underline{\mathrm{Pic}}^\tau_{X/S}$ est-il
universellement ouvert sur $S$ ?'' occurs on 57:309--311. The following
example disproves it in the unrestricted proper-flat setting, already in
characteristic zero. It also disproves unrestricted constancy of
$\dim\operatorname{Pic}^0_{X_s}$ in 57:313(i). It has nonreduced fibres, so
the separately written normal-morphism version remains a distinct question.

Let $E/\mathbf C$ be the elliptic curve $y^2z=x^3-xz^2$, with origin $o$, let
$S=E$, and let
$\mathcal P=\mathcal O_{E\times S}(\Delta-E\times o-o\times S)$ be the
normalized Poincar\'e bundle. On $V=E\times\mathbf P^1\times S$ put
\[
 M=\operatorname{pr}_{E,S}^*\mathcal P\otimes
       \operatorname{pr}_{\mathbf P^1}^*\mathcal O_{\mathbf P^1}(-2),
 \qquad \mathcal X=\operatorname{Spec}_V(\mathcal O_V\oplus M),\quad M^2=0.
\]
The split square-zero algebra is finite locally free of rank two over $V$.
Thus $\mathcal X/S$ is projective and flat. The vanishing
$H^0(\mathbf P^1,\mathcal O(-2))=0$ gives
$f_*\mathcal O_{\mathcal X}=\mathcal O_S$ after every base change.
Over a geometric point $s$, the splitting of units
$\mathcal O_{\mathcal X_s}^{\times}=\mathcal O_{V_s}^{\times}\times(1+M_s)$
gives
\[
 \operatorname{Pic}_{\mathcal X_s}
 \simeq\operatorname{Pic}_{E\times\mathbf P^1}\times H^1(V_s,M_s)
\]
as group schemes, with the last factor additive. A degree-zero line bundle
on $E$ has both $H^0$ and $H^1$ zero unless trivial, and both one-dimensional
when trivial. K\"unneth therefore gives
\[
 (h^0(M_s),h^1(M_s),h^2(M_s))=
 \begin{cases}(0,0,0),&s\ne o,\\(0,1,1),&s=o.\end{cases}
\]
Consequently $\dim\operatorname{Pic}^0$ is one away from $o$ and two at $o$.
SGA~3 VI$_B$, Proposition~4.2, says that a group scheme locally of finite
presentation is universally open along the identity precisely when its fibre
dimension is locally constant there. Thus the tau-part is not universally
open. The same statement can be seen in the extra Picard factor over the
local DVR at $o$: its equation is $tz=0$ in $\mathbf G_{a,R}$, and the open
$z\ne0$ maps to the closed point. The dimension argument alone suffices and
does not require this local equation.

There is a positive statement where the manuscript specifies normal fibres:
when the fibrewise Picard identity components have no additive subgroup,
FGA VI Corollary~2.7 makes their dimension locally constant, and SGA~3
VI$_B$~4.2 gives universal openness along the neutral part. Multiplication by
an integer invertible on the base is \'etale, so the same holds along components
whose order in the component group is prime to the residue characteristic.
In characteristic zero this covers the whole tau-part. For normal fibres
in characteristic $p$, the question remains open here on the components
whose classes have $p$-power order.

\subsection{E4: the product formula, with the correct hypotheses}

For $X,Y$ proper flat and finitely presented over $S$, each with universal
$H^0=\mathcal O_S$, the canonical map
\[
 \operatorname{Pic}^\tau_{X/S}\times_S\operatorname{Pic}^\tau_{Y/S}
 \xrightarrow{\sim}\operatorname{Pic}^\tau_{X\times_SY/S}
\]
is an isomorphism of fppf sheaves. Here is the argument completing the
question on 57:309(g). E2 represents the correspondence quotient by a
separated unramified group scheme. Its geometric fibres are torsion-free.
Indeed over an algebraically closed field, choose base points and normalize
a correspondence as in E2. An $n$-torsion class gives a based map
$a:X\to\operatorname{Pic}^0_Y[n]$. This target is affine. To see this, apply
Chevalley's theorem to the reduced connected Picard group: its affine part
maps to a finite $n$-torsion subscheme of its abelian quotient, and a nilpotent
thickening of an affine scheme is affine. Since $X$ is proper with
$H^0(X,\mathcal O_X)=k$, every map from $X$ to an affine scheme is constant.
The based map is therefore zero.

An algebraically trivial Picard class maps to zero in the discrete
correspondence scheme, and a tau-class maps to torsion, hence to zero by
the preceding paragraph. Over any base, its correspondence section is zero
on all geometric fibres. The equalizer with the identity is open because
the correspondence scheme is unramified, and contains every point, so the
section is zero. Fppf locally one can choose base points, and the normalized
bundle is then an external product. Restriction to the two axes says that
both factors are tau-trivial; uniqueness and fppf descent finish the proof.

Geometric connectedness alone, one of the uncertain marginal readings, is
insufficient. Take $X=\operatorname{Spec}k[\epsilon]/(\epsilon^2)$ and an
elliptic curve $Y/k$. Then $\operatorname{Pic}^\tau_X=0$, while the Picard
functor of $X\times Y$ has an extra additive factor $H^1(Y,\epsilon\mathcal O_Y)$
of dimension one. Its neutral dimension is two, whereas that of
$\operatorname{Pic}^\tau_Y$ is one. The universal $H^0=\mathcal O$ condition,
or geometric reducedness together with connectedness, removes this defect.

\subsection{E4: a proper normal family with nonproper tau-part}

\paragraph{Theorem E4.N (explicit counterexample).}
There is a projective flat morphism $X\to\operatorname{Spec}\mathbf C[[t]]$
with geometrically integral normal fibres and a section, for which the
represented separated finite-type group scheme $\operatorname{Pic}^\tau_{X/R}$
is not proper. In particular the normal-fibre branch of 57:310(5) is false,
and normal projectivity of $X/R$ does not imply projectivity of its tau-part.
Since the base has characteristic zero, its tau-part is universally open;
thus the implication proposed in 57:310(5') is false as well.

\paragraph{Construction.}
Put $R=\mathbf C[[t]]$, $K=\mathbf C((t))$. In $\mathbf P^5_R$ take the
three quadrics $\sum_i a_{ji}x_i^2=0$ with coefficient matrix
\[
 A(t)=\begin{pmatrix}
 1&1&1&1&1&1\\
 0&1&2&3&5&7\\
 0&1&2+t&9&25&49
 \end{pmatrix}.
\]
Call their intersection $Y$, and let $g$ fix the first three homogeneous
coordinates and negate the last three. Every three-column determinant of
$A(t)$ is nonzero in $K$; at $t=0$ exactly the determinant of columns
$0,1,2$ vanishes, and it equals $t$. These twenty exact determinants are
listed in the accompanying computation output. Any two columns are
independent. The generic surface is therefore a smooth complete intersection
of three quadrics: a deficient Jacobian would be supported on at most two
coordinates, and then the equations force all coordinates to vanish.
The same argument on the special fibre shows that its only singular points
are the four points supported on $0,1,2$ with squared coordinates
\[
 (x_0^2,x_1^2,x_2^2)=(1,-2,1).
\]
Indeed the kernel of the first three columns at zero is spanned by
$(1,-2,1)$, whose three entries are nonzero. After projectivization there
are precisely four choices of signs.

On the affine chart $x_2=1$ at each of these points, the first two equations
solve for $x_0,x_1$ in the completed local ring: their $2\times2$ Jacobian
has nonzero determinant since $x_0x_1\ne0$. Subtracting the second quadric
from the third leaves exactly
\[
 t+6x_3^2+20x_4^2+42x_5^2=0.                         \tag{E4.N.1}
\]
Thus $Y$ is regular at each of these points and the special fibre has an
ordinary double point there. Away from them, fibre smoothness makes the
total space regular. The squaring map $[x_i]\mapsto[x_i^2]$ on $\mathbf P^5_R$ is finite
locally free. The three linear equations with matrix $A(t)$ cut out a
$\mathbf P^2_R$, because the last three columns have unit determinant.
Thus $Y$, their inverse image under squaring, is finite flat over that
$\mathbf P^2_R$ and is flat over $R$; it is also a relative complete
intersection.
The fibres are connected by their Koszul resolution, are Cohen--Macaulay,
and are regular in codimension one; consequently they are normal and
geometrically integral. The same resolution gives $H^0=\mathbf C$ and
$H^1(\mathcal O)=0$ on both geometric fibres. Adjunction gives trivial
canonical bundle. In particular the generic fibre is a K3 surface.

The fixed locus of $g$ in projective space consists of the two coordinate
planes. The corresponding three-column matrices show that $g$ has no
fixed point on $Y_K$ and exactly the four displayed fixed points on $Y_0$.
The finite quotient $q:Y\to X=Y/\langle g\rangle$ exists as a projective
scheme: the linear action supplies an invariant Veronese homogeneous
coordinate ring. Since $2$ is invertible, invariants commute with every
base change. The quotient is flat over $R$ (its local invariant rings are
$R$-torsion-free) and its geometric fibres are normal integral schemes,
because invariant subrings of normal domains are normal. Reynolds
averaging and the preceding Koszul calculation give universal
$f_*\mathcal O_X=\mathcal O_R$ and $H^1(X_s,\mathcal O)=0$.
Choose a smooth closed point of $X_0$ away from the four quotient singularities.
Henselian lifting through the smooth locus gives an $R$-section of $X$.
This ensures that every $R$-point of the relative Picard sheaf is represented
by an actual line bundle.

\paragraph{Failure of extension.}
The connected finite \'etale double cover $q_K:Y_K\to X_K$ has a nontrivial
anti-character line bundle $T_K$, with $T_K^{\otimes2}\simeq\mathcal O$.
It is nontrivial even geometrically: otherwise the geometrically connected
K3 cover would split after extending $K$. Suppose it extended to a line
bundle $L$ on $X$. Then $q^*L$ is trivial on $Y_K$. Since $Y$ is regular,
a rational trivialization gives a Cartier divisor supported on $Y_0$.
There is only one reduced irreducible vertical divisor, $Y_0$, and it is
$\operatorname{div}(t)$. Thus $q^*L\simeq\mathcal O_Y$.

The natural $g$-linearization on a pullback line bundle acts trivially on
its fibre at any fixed point of $g$. On the other hand, any linearization
on $\mathcal O_Y$ is given by a unit $c\in H^0(Y,\mathcal O_Y)^*=R^*$ with
$c^2=1$; hence $c=+1$ or $-1$. Its generic restriction must be the sign
linearization defining $T_K$, so $c=-1$. This contradicts triviality on
fibres at the four fixed points. Therefore $T_K$ does not extend to $X$.

The projective geometrically integral family has a represented separated
Picard scheme, and its tau-part is open of finite presentation by the
standard finiteness theorem. Since $T_K$ is a tau-point and cannot extend,
the tau-part fails the valuative criterion for properness. There is no
conflict with E1: $T_K$ is a nonzero N\'eron--Severi torsion class; it is not
in the neutral component. Indeed $H^1(X_K,\mathcal O)=0$ and characteristic
zero make the neutral component trivial. The universally open assertion
follows from FGA VI Theorem~2.5 and its corollary: normal geometric fibres
have proper neutral Picard varieties, and every component order is
invertible over this base.

\subsection{E4: dimension jump even with geometrically reduced fibres}

\paragraph{Theorem E4.R.}
The constancy question on 57:313(i) is false even with geometrically reduced
fibres, the marginal hypothesis ``du moins pour fibres s\'eparables''.
Consequently the universal-openness questions 57:309(c), 310(3), and 311
also have negative answers under this weaker hypothesis than normality.

Keep the elliptic curve $E$, base $S=E$, origin $o$, and normalized
Poincar\'e bundle $\mathcal P$ on $C=E\times S$ used above. Let $V\to C$ be
the projective-line bundle compactifying the total space of
$\mathcal P^\vee$, and write $C\subset V$ for its zero section. Its conormal
bundle is $\mathcal P$. The doubled zero section $D=2C$ is the split ribbon
\[
 \mathcal O_D=\mathcal O_C\oplus\mathcal P,\qquad\mathcal P^2=0,
\]
with the splitting induced by $V\to C$. Define a finite $V$-scheme
\[
 X=\operatorname{Spec}_V\!\left(
 \mathcal O_V\times_{\mathcal O_D}\mathcal O_V\right).       \tag{E4.R.1}
\]
Writing $I=\mathcal O_V(-D)$, the fibre-product algebra is
$\mathcal O_V\oplus I$ as an $\mathcal O_V$-module. Locally, if $d$ is an
equation for $D$, its equation is $z(z-d)=0$. In particular $X\to V$ is
finite locally free of rank two, so $X/S$ is projective and flat.

Formation of this algebra commutes with arbitrary base change on $S$:
in the exact sequence
\[
 0\longrightarrow I\longrightarrow\mathcal O_V
     \longrightarrow\mathcal O_D\longrightarrow0
\]
all three terms are flat over $S$; $D$ is a square-zero extension of the
smooth $C/S$ by an invertible module. Thus $I_T$ remains the kernel after
any base change $T\to S$, and the fibre-product algebra remains the same
fibre product. Each geometric fibre of $V$ is integral and $d$ is a
nonzerodivisor. Hence the algebra in (E4.R.1), on that fibre, embeds in
$\mathcal O_{V_s}\times\mathcal O_{V_s}$ and is reduced. The two components
meet along the nonempty ribbon $D_s$, so the fibre is connected. Moreover
\[
 0\longrightarrow\mathcal O_X\longrightarrow
 \mathcal O_V\oplus\mathcal O_V\longrightarrow\mathcal O_D\longrightarrow0
\]
shows $H^0(X_T,\mathcal O)=H^0(T,\mathcal O_T)$ for every $T$: the map on
the two constant summands is their difference in the copy of
$\mathcal O_T\subset H^0(D_T,\mathcal O)$, so its kernel is the diagonal.
This proves universal cohomological flatness in degree zero directly.

For a geometric point $s$, the units Mayer--Vietoris sequence gives
\[
 H^0(V_s,\mathcal O)^*\times H^0(V_s,\mathcal O)^*
 \longrightarrow H^0(D_s,\mathcal O)^*
 \longrightarrow\operatorname{Pic}(X_s)
 \longrightarrow\operatorname{Pic}(V_s)^2
 \longrightarrow\operatorname{Pic}(D_s).
\]
Here $H^0(D_s,\mathcal O)^*=k^*\times H^0(E,\mathcal P_s)$ with the
second factor additive; the first map has image exactly the scalar $k^*$.
Also $\operatorname{Pic}^0(V_s)=\operatorname{Pic}^0(E)$, and restriction to
$D_s$ is the split pullback from $E$. The kernel on neutral components of
the last map is therefore the diagonal elliptic curve. Pullback along the
folding map $X_s\to V_s\to E$ splits this diagonal, giving an isomorphism
of group schemes
\[
 \operatorname{Pic}^0_{X_s/k}
  \simeq E\times\mathbf G_a^{\,h^0(E,\mathcal P_s)}.
\]
The units sequence and the splittings used here are functorial after a
$k$-algebra extension, so this is a group-scheme computation, not only a
calculation of rational points. Since $\mathcal P_s$ is a nontrivial
degree-zero bundle for $s\ne o$, the exponent is zero there and is one
at $o$. The neutral dimension jumps from one to two. SGA~3 VI$_B$~4.2
then rules out universal openness at the identity. Although the fibres
are reduced, they are not normal: the two smooth components have been
glued along a doubled divisor. This is the precise scope of the example.

\subsection{E4: the abelian and additive rank inequalities}

FGA VI Remark~2.8, pp.234--235, fixes the convention in 57:313(ii):
unprimed invariants belong to the special point $s$, primed invariants to
its generization $\eta$. Thus the manuscript asks for
$\alpha_s\le\alpha_\eta$ and $\lambda_s\ge\lambda_\eta$.
Here $\alpha,\mu,\lambda$ are the dimensions of the abelian quotient,
torus, and unipotent part of the reduced connected geometric Picard group.
Both inequalities hold in the represented proper-flat Picard setting;
in particular they hold under the geometrically reduced hypothesis in
the margin. They do not imply constancy of their sum, as E4.R shows.

We use Theorems~1.1 and~4.1 of Gille--Guralnick
(arXiv:2211.09876, pp.~1, 7--8), \emph{Unipotent Subgroups of
Stabilizers}. For separated group schemes of finite presentation,
they prove upper semicontinuity of two dimensions: the maximal
unipotent-subgroup dimension $d_u$, and the maximal dimension $d_c$
of an affine commutative subgroup. No flatness hypothesis is required. For a commutative algebraic group,
\[
 d_u=\lambda,\qquad d_c=\mu+\lambda,\qquad
 \alpha=\dim G-d_c.                                  \tag{E4.rank.1}
\]
The elementary Chevalley decomposition proves these identifications.
The cited proof takes the closure of a generic subgroup over a trait,
uses Raynaud's affineness theorem (proved in Prasad--Yu, Proposition~3.1,
p.5), and the finite smooth normalization after a finite extension
(Conrad's Appendix, Corollary~A.4, p.29). Hence no unproved extension
of SGA~3's smooth-group announcement is being used.

First work over a strictly henselian trait and suppose $P^\tau$ is
separated; it is of finite presentation. The $d_u$ theorem immediately
gives $\lambda_s\ge\lambda_\eta$. Let $G$ be the schematic closure of
$P^\tau_\eta$ in $P^\tau$. This is a closed finite-type flat subgroup:
flatness is torsion-freeness over the DVR, and multiplication and inversion
preserve the closure because products of flat schemes have schematically
dense generic fibre. Put $B=(P_s)^0_{\mathrm{red}}$ and
$H=(G_s)^0_{\mathrm{red}}$, over an algebraic closure of the residue field;
these are smooth connected groups.

The quotient $B/H$ is unipotent.
Here is a direct verification of the relevant assertion, also recorded
in Raynaud 1970, Proposition~6.3.1, p.52. A line bundle of order prime to
the residue characteristic on $X_s$ is represented by a finite \'etale
Kummer torsor. The equivalence between finite \'etale covers of a proper
scheme over a henselian local base and covers of its closed fibre
(Artin 1969, Theorem~3.1, p.30) lifts that torsor, with its group action,
to $X$. It therefore lifts the torsion Picard point. Its generic value is
in $P^\tau_\eta$, so its section lies in $G$. For every sufficiently large
prime $\ell$ invertible in the residue field, the finite group of
components of $G_s$ has no $\ell$-torsion, and hence
$(P_s)^0[\ell]\subset H$. Multiplication by $\ell$ is surjective on the
connected algebraic group $H$, so the quotient has no $\ell$-torsion.
A positive-dimensional torus or abelian quotient would have nonzero
$\ell$-torsion for every such $\ell$. The quotient is therefore unipotent.
It follows that $\alpha(P_s)=\alpha(G_s)$.

The flat group $G$ has constant fibre dimension. Apply upper
semicontinuity of $d_c$ and subtract in (E4.rank.1):
\[
 \alpha_s=\alpha(G_s)
 =\dim G_\eta-d_c(G_s)
 \le\dim G_\eta-d_c(G_\eta)=\alpha_\eta.
\]
This proves the two claimed inequalities over traits; constructibility
and the usual trait reduction give the specialization statement over a
locally noetherian base.

For completeness, separation of the whole Picard scheme is unnecessary.
If the proper-flat Picard functor $P$ is represented over the trait,
Raynaud's Proposition~5.2 makes the closure $E$ of the generic identity
\'etale, and his \S4.1 constructs a separated group scheme $Q=P/E$.
Here the passage from a separated group algebraic space to a scheme uses
precisely Raynaud Theorem~3.3.1 over the trait.
The image $Q'$ of the finite-presentation open subgroup $P^\tau$ is a
quasi-compact open subgroup of $Q$, hence a separated finite-presentation
group scheme. On every geometric fibre the kernel is zero-dimensional;
thus $P^\tau\to Q'$ preserves $\alpha,\mu,\lambda$ on the reduced connected
groups. Apply the preceding argument to $Q'$; prime-to-characteristic
torsion lifts after choosing $\ell$ also prime to the finite special
kernel. This proves the same conclusion when the represented tau-part
is nonseparated. The statement is not extended here to an arbitrary
nonrepresentable Picard functor without checking the same quotient data.

\paragraph{The prime-to-characteristic rank inequality.}
The original printed formulas in FGA VI Remark~2.8, pp.234--235,
are unambiguous: with $s$ special and $\eta$ generic,
\[
 2\alpha_s+\mu_s\le 2\alpha_\eta+\mu_\eta.       \tag{E4.rank.2}
\]
In the geometrically reduced setting of 57:313, this follows directly
from separation of $P^\tau$ proved above. Over a strictly henselian
trait, choose a prime $\ell$ invertible on the trait and not dividing
the orders of the torsion in either geometric N\'eron--Severi group.
The kernel $P[\ell]$ is \'etale by FGA VI Theorem~2.5 and is separated
because it lies in $P^\tau$. Each special point therefore extends to
a section, and distinct sections have distinct generic values.
The geometric cardinalities are respectively
$\ell^{2\alpha_s+\mu_s}$ and
$\ell^{2\alpha_\eta+\mu_\eta}$: prime-to-characteristic
multiplication on the unipotent part is an isomorphism, and on the
torus and abelian variety it has these familiar orders.
This proves (E4.rank.2), including its direction.
More generally the same proof applies whenever the relevant
prime-to-characteristic torsion subspace is separated.
This assertion is distinct from constancy of Picard dimension,
which E4.R refutes even with geometrically reduced fibres.

\subsection{Flatness and openness: a Heisenberg degeneration}

\paragraph{Theorem E4.H.}
There exists a smooth projective family of geometrically connected surfaces
$X/R$, with $R=W(k)$ and $k$ an algebraically closed field of characteristic
$2$, such that
\[
 \operatorname{Pic}^{\tau}_{X/R}
 \simeq\mu_{2,R}^{\,2}\times_R K,
 \qquad K=\operatorname{Spec}R\amalg\operatorname{Spec}k.       \tag{E4.H.1}
\]
The second summand of $K$ has the group law inherited from the nonzero
component of the constant group $(\mathbf Z/2)_R$; its product with itself
lands in the identity summand. Consequently $\operatorname{Pic}^{\tau}_{X/R}$
is not open over $R$, even though $X/R$ is smooth. Over
$A=R/4R$, its fibrewise neutral subgroup is represented by
$\mu_{2,A}^{\,2}$ and is finite flat, whereas its tau-part is not flat.
This refutes both the normal-family universal-openness question on 57:310(3)
and the implication written on 57:311, even with represented neutral part.

\paragraph{The finite group and its characters.}
Define a finite flat group $G/R$ by
\[
 G=\mu_{2,R}\times (\mathbf Z/2)_R^2,
\quad
 (z,a,b)(z',a',b')=
 \bigl(zz'(-1)^{ba'},a+a',b+b'\bigr).                        \tag{E4.H.2}
\]
Here the displayed product decomposition is a decomposition of schemes;
the group is not commutative. The exponent is taken modulo two.
The elementary bilinear cocycle identity proves associativity. Its central
$\mu_2$ has quotient $(\mathbf Z/2)^2$. The generic geometric group is the
dihedral group of order eight, with commutator subgroup its central order-two
subgroup. The special group is the commutative group
$\mu_{2,k}\times(\mathbf Z/2)_k^2$, since $-1=1$ in $k$.
This is the case $p=2$, $e=p-1$ of Raynaud's construction
\cite[Lemma~4.2.1, pp.~135--136]{BCERaynaud1979}.

A character of $G_T$ restricts on the central $\mu_{2,T}$ to the locally
constant exponent $n\in\mathbf Z/2$. Its values on $(1,1,0)$ and
$(1,0,1)$ are arbitrary elements $u,v\in\mu_2(T)$. The commutator relation
imposes exactly $(-1)^n=1$: the component $n=0$ is defined over all of $R$,
and $n=1$ exists precisely over $R/2R=k$. There are no further relations.
Thus the character functor, including nonflat tests, is
\[
 \underline{\operatorname{Hom}}(G,\mathbf G_m)
 =\mu_{2,R}^{\,2}\times_R K.                               \tag{E4.H.3}
\]
In particular its coordinate module is $R^4\oplus k^4$.
The code \texttt{check\_raynaud\_group.py} checks all 64 cocycle identities,
all 512 associativity instances of the generic eight-element group, its
nontrivial commutator, and its four characters. The Hopf-scheme calculation
(E4.H.3), rather than point counting, supplies the special length eight.

\paragraph{A completely specified complete-intersection realization.}
The following finite formula fixes the degree, every section, and every
coefficient. Use the 24 coordinates $x_{r,n,a,b}$, where $1\leq r\leq3$
and $n,a,b\in\{0,1\}$. The central $\mu_2$ has weight $n$, and
\[
 A(x_{r,n,a,b})=x_{r,n,a+1,b},\qquad
 B(x_{r,n,a,b})=(-1)^{na}x_{r,n,a,b+1}.                    \tag{E4.H.4}
\]
The indices $a,b$ are modulo two. These are three copies of the regular
coordinate representation, so $P=\mathbf P^{23}$ is the projectivization
of the threefold direct sum of the contragredient regular representation
in Raynaud's Lemma~4.2.2. Its free-action open $U$ has complement of
fibrewise codimension at least three.

Order the coordinates lexicographically. Let $\mathcal A$ be the ordered
finite set of degree-192 monomials with even total $\mu_2$ weight. Form
\[
 F=\sum_{\alpha\in\mathcal A}u_\alpha x^\alpha,
 \qquad F\,A(F)\,B(F)\,AB(F)=\sum_{|\beta|=4}u^\beta p_\beta(x).
                                                               \tag{E4.H.5}
\]
The actions fix the $u$ variables. On even-weight forms $A,B$ commute,
so every integral polynomial $p_\beta$ is $G$-invariant of degree 768.
Retain zero polynomials and repetitions; no basis choice is involved.
Let $\mathcal B$ be the ordered set of the multiindices $\beta$ above.
Define
\[
 k=\overline{\mathbf F_2(c_{i\beta}:1\leq i\leq21,
                                  \ \beta\in\mathcal B)},
 \quad R=W(k),\qquad
 s_i=\sum_{\beta\in\mathcal B}[c_{i\beta}]p_\beta.          \tag{E4.H.6}
\]
Here all variables are algebraically independent and brackets are
Teichm\"uller representatives. These are finite definitions, with
\[
 h=|\mathcal A|=\tfrac12\{\binom{215}{23}+\binom{107}{11}\},
 \qquad |\mathcal B|=\binom{h+3}{4}.
\]

To type the descent precisely, put $M=P/\mu_2$ and $Q=M/H=P/G$, where
$H=(\mathbf Z/2)^2$. The central norm of $\mathcal O_P(1)$ is a bundle
$\mathcal N$ on $M$ with pullback $\mathcal O_P(2)$, naturally
$H$-linearized. Put $\mathcal L=\operatorname{Norm}_H(\mathcal N)$ on
$Q$. It is ample with pullback $\mathcal O_P(8)$. These are Raynaud's
norm bundles with their factorization made explicit
\cite[proof of Proposition~4.2.3, p.~137]{BCERaynaud1979}.
Even-weight degree-192 forms are sections of $\mathcal N^{96}$.
Their $H$-norms, hence all $p_\beta$, are sections of $\mathcal L^{96}$.
These integral formulas avoid any assumption about reduction of a chosen
basis of sections. Define
\[
 X=V(s_1,\ldots,s_{21})\subset Q,
 \qquad Y=q^{-1}(X)\subset P.                              \tag{E4.H.7}
\]
Thus $Y$ is specified by 21 equations of the fixed degree 768.

\paragraph{Why this fixed degree suffices.}
Write $V=U/G$. The maps $U\to U/\mu_2\to V$ are respectively a degree-two
$\mu_2$-torsor and a degree-four finite \emph{\'etale} $H$-torsor.
The scheme $V$ is smooth of relative dimension 23: flatness descends
from $U$, and geometric regularity descends through the faithfully flat
finite maps on fibres. The complement $Q_k\setminus V_k$ has dimension
at most 20.

The sections $p_\beta$ generate $\mathcal L^{96}$ on $Q_k$. Indeed even
coordinate powers generate $\mathcal N^{96}$ on $M_k$. Since $k$ is
infinite, a linear combination can be chosen nonzero at all points of
any finite $H$-orbit. Its norm is nonzero at the corresponding point of
$Q_k$, and every norm is a linear combination of the coefficient
sections $p_\beta$.

They also span every first jet at $v\in V(k)$. The first infinitesimal
neighborhood of $v$ has length 24; its inverse image $Z$ in $P_k$ has
length $8\cdot24=192$. A length-$\ell$ subscheme of projective space
over an infinite field is interpolated by $\mathcal O(d)$ for
$d\geq\ell-1$: choose a hyperplane avoiding its support and let $W_d$
be the span of degree-at-most-$d$ monomials in its length-$\ell$
coordinate algebra. If $W_{d+1}=W_d$, this space is stable under every
coordinate and is the whole algebra. Otherwise its dimension strictly
increases. Since $\dim W_0=1$, it reaches the algebra by $d=\ell-1$.
Thus degree-192 forms interpolate all sections on $Z$.

Taking $\mu_2$ invariants preserves that surjection: projection onto
weight zero is exact even in characteristic two. Sections of
$\mathcal N^{96}$ therefore interpolate arbitrary first jets at the
four points above $v$ in $U/\mu_2$. Prescribe the desired jet at one
point and unit jets at the other three, in $H$-compatible
trivializations. The norm has the desired jet at $v$. Consequently
the span of the $p_\beta$ surjects onto every first jet on $V_k$.

The incidence proof of Bertini now makes 21 generic sections smooth on
$V_k$, while global generation and the dimension bound 20 make them
avoid its complement. The independent coefficients in (E4.H.6) specify
such a generic tuple. Amplitude makes the intersection a nonempty
surface. Thus $X_k$ is smooth and contained in $V_k$. Since $Q$ is smooth
near $X_k$ and the 21 equations cut its smooth special fibre in codimension
21, they form a fibrewise regular sequence there. The local flatness
criterion makes $X/R$ flat near $X_k$, and the smooth special fibre then
makes $X/R$ smooth on that neighborhood. These properties hold throughout
$X$: their bad loci are closed in the projective
$R$-scheme $X$, and a nonempty closed image in the local base must meet
the closed point. Therefore $X/R$ is smooth projective, $Y\to X$ is a
$G$-torsor, and the 21 pullback equations are a relative complete
intersection in $P^{23}_R$. This verifies all hypotheses of Raynaud's
Proposition~4.2.3 with an explicit power $m=96$.

A complete intersection of dimension at least two has
$H^0(\mathcal O)=k$, $H^1(\mathcal O)=0$, and trivial Picard tau-group,
including when singular or nonreduced; the last statement is
\cite[Theorem~1.1.8(iii), proved in \S3.4]{BCECotnerZavyalov2024}.
Hence $Y/R$ has universal $H^0=\mathcal O_R$ and
$\operatorname{Pic}^\tau_{Y/R}=0$, by the fibrewise assertion and the
unramified identity-section criterion. Geometric connectedness of $Y$
and its surjection onto $X$ make the smooth fibres of $X$ geometrically
integral and normal. The coordinate-check script verifies the signed
representation and the finite index cardinalities. It does not expand
the enormous norm polynomial or claim a computational smoothness test;
the interpolation and incidence arguments supply that proof.

\paragraph{Picard descent and the two failures.}
The kernel of
$\operatorname{Pic}_{X/R}\to\operatorname{Pic}_{Y/R}$ is the character
functor of $G$. Indeed a line bundle trivialized after pullback to $Y$ is a
$G$-linearization of $\mathcal O_Y$; universal $H^0(Y,\mathcal O_Y)=R$
identifies the linearization with a unit on $G$ satisfying exactly the
character identity. This argument is fppf-local on every test scheme and
therefore identifies the full sheaves, including their nilpotents.
Since $\operatorname{Pic}^\tau_{Y/R}=0$, every tau-class lies in this
kernel. Conversely the character group in (E4.H.3) is killed by two and
hence lies in tau. This proves (E4.H.1) and its compatibility with arbitrary
base change.

The open-and-closed component
$\mu_{2,k}^{\,2}\subset\operatorname{Pic}^\tau_{X/R}$ has image the closed
point of the trait. The map is therefore not open, and a fortiori not
universally open. This closes the positive-residue-characteristic gate;
no implication from nonflatness alone is being used.
After base change to $A=R/4R$, the neutral component of the only geometric
fibre lies in the first summand, so the fibrewise neutral subfunctor is
exactly $\mu_{2,A}^{\,2}$, finite locally free of rank four. The other
summand has coordinate module $k^4$ and is not $A$-flat: a nonzero finite flat module over the
local ring $A$ is free, whereas this module is annihilated by $2\ne0$. Thus flatness of the represented
neutral part does not imply flatness of tau. The construction depends on
Raynaud's noncommutative finite-flat group, not on the disputed tangent-space
argument for ordinary Igusa deformations.

\subsection{The quotient question in its written generality}

\paragraph{Theorem E4.Q (refutation of 57:309(b)).}
The manuscript asks for ``existence du pr\'esch\'ema quotient''
$\underline{\mathrm{Pic}}_{X/S}/\underline{\mathrm{Pic}}^\tau_{X/S}$.
In this generality the answer is negative, even when the tau-part is flat.
Let
\[
 S=\operatorname{Spec}\mathbf R[[t]],\qquad
 X=V(x^2+y^2-tz^2)\subset\mathbf P^2_S .                  \tag{E4.Q.1}
\]
This is Mumford's example as proved in
\cite[Example~4.14, pp.~31--33]{Kleiman2005}. Its Picard fppf sheaf
$P$ is an algebraic space, not a scheme. After the finite \'etale base change
$S'=\operatorname{Spec}\mathbf C[[t]]$, its component of total degree $n$
is obtained by gluing copies $S'_{a,b}$, $a+b=n$, along their generic point.
Complex conjugation interchanges the origins $(a,b)$ and $(b,a)$.
For $a\ne b$ these origins have no common affine neighborhood: their two
trait sections agree generically, contrary to separatedness of such an
affine neighborhood. If $P$ were a scheme, an affine neighborhood of their
common descended point would pull back to just such an affine neighborhood.

In this example $P^\tau$ is exactly the identity subgroup $S$. Indeed the
geometric generic conic has Picard group $\mathbf Z$, and the geometric
special fibre, two projective lines meeting in one point, has Picard group
$\mathbf Z^2$ by multidegree. Both have trivial neutral and torsion parts.
In the displayed scheme $P_{S'}$, the tau-subfunctor is consequently the
identity copy $S'_{0,0}$; the explicit representation also rules out any
additional infinitesimal structure. This identification descends to $S$.
Therefore, as fppf sheaves of groups,
\[
 P/P^\tau=P/S\simeq P,
\]
which is not a scheme. Here $S$ denotes the trivial group object, not a
nontrivial quotient of the underlying base. Its identity immersion into
$P$ is open and not closed. Thus this example is consistent with the flat-only algebraic-space assertion
in the prompt summary of H6. The stronger scheme conclusion in the fuller
\texttt{HUNCHES.md} formulation additionally assumes closedness, which fails here.

The stronger variant that presupposes $P$ itself schematic, or imposes
closedness of $P^\tau$ while dropping flatness, is a distinct question.
No answer to those extra variants is inferred from (E4.Q.1). In particular,
this example has a geometrically reducible special fibre and gives no
counterexample to the geometrically integral or normal question on 57:315.

\paragraph{Two further consequences of E4.H.}
The character group in (E4.H.1) is killed by two. Its identity section is
not open: in the component $\mu_{2,R}^{\,2}$, the four geometric generic
points specialize to the same special identity. Thus the image of the
whole group under $[2]$ is not open. Question 57:310(4) is therefore false
even for a smooth projective family, extending the simpler cuspidal
field example.

For the margin question about ``Multiplicit\'e de
$\underline{\mathrm{Pic}}^\tau$'' and specialization on 57:310, the same
example makes several precise natural interpretations explicit. The
generic fibre has four reduced points, each of multiplicity one. The
special fibre consists of two copies of $\mu_{2,k}^{\,2}$, each of length
four. Total length changes from four to eight, the number of geometric
components from four to two, and the length of the neutral geometric
component from one to four. Thus none of these numerical multiplicities
is specialization invariant without further hypotheses. The manuscript's
additional words about existence and descent do not specify an object;
we do not replace them by a different theorem.

\subsection{E4: closedness of the numerical locus without projectivity}

\paragraph{Theorem E4.T.}
Let $f:X\to S$ be a proper flat morphism of finite presentation with
geometrically integral fibres. The base scheme $S$ is arbitrary. For every
invertible sheaf $L$ on $X$, the set
\[
 S^\tau(L)=\{s\in S:L_{\bar s}\text{ is numerically trivial}\}
                                                        \tag{E4.T.1}
\]
is open and closed. Equivalently, the open subgroup algebraic space
$\operatorname{Pic}^{\tau}_{X/S}$ is closed in
$\operatorname{Pic}_{X/S}$. Thus projectivity is unnecessary in the
question ``projectif est-il n\'ecessaire ?'' on 57:314. Here numerical
triviality and tau-equivalence agree for proper schemes over a field;
see \cite[Remark~6.14, pp.~57--58]{Kleiman2005}.

\paragraph{1. The trait assertion.}
Let $R$ be a discrete valuation ring, with fraction field $K$, and let
$X/R$ satisfy the hypotheses. Suppose $L_K$ is numerically trivial.
We prove that $L_k$ is numerically trivial. Replacing $R$ by its completion
preserves the special fibre and the hypotheses; numerical triviality of
the generic bundle is preserved under field extension. We may therefore
assume $R$ complete, hence excellent. Flatness and integrality of the generic
fibre make $X$ integral. If $\xi$ is the generic point of its reduced
integral special fibre, $\mathcal O_{X,\xi}$ is a one-dimensional local
domain with maximal ideal generated by a uniformizer of $R$. It is a DVR.

Zheng's Galois alteration lemma
\cite[Lemma~3.8, pp.~313--314]{BCEZheng2009} gives a finite extension of
complete traits $R\subset R'$, a regular integral scheme $Y$ projective
over $R'$, and a proper generically finite map
\[
 \pi:Y\longrightarrow X
\]
equivariant for a finite group $G$ acting trivially on $X$. Moreover
$k(Y)^G/k(X)$ is purely inseparable. This is precisely the definition of
Galois alteration used on p.~313 of that source. To apply the lemma take
the input group trivial and the dense open $X_K$. Its compactification
$Y$ is the projective end of a plurinodal fibration over $R'$. The source
$X'$ of the open immersion in the general statement is proper over $X$,
hence over $R$. Since $R'/R$ is separated, $X'$ is proper over $R'$ as
well. Its map into the separated projective end is therefore proper.
The dense open immersion is consequently an isomorphism. The lemma also makes the
boundary a strict normal-crossings divisor. Only regularity of $Y$ and
the intersection theory below are needed: regularity of a geometric
special fibre is not assumed.

Put $M=\pi^*L$ and write
\[
 Y_{k'}=\sum_{i=1}^r m_iY_i,\qquad m_i>0.                \tag{E4.T.2}
\]
The special fibre is connected. Indeed $H^0(Y,\mathcal O_Y)$ is finite
over the complete local ring $R'$ and is a domain, hence a local ring;
Stein factorization then gives connectedness. Each $Y_i$ is a Cartier
divisor on the regular $Y$. Let $d=\dim Y_{K'}$. The case $d=0$ is
immediate, since every line bundle on a proper zero-dimensional scheme
is numerically trivial. Assume $d\geq1$, and choose a relatively ample
line bundle $H$ on $Y$.

\paragraph{2. The vertical correction.}
Define the rational symmetric matrix $A$ and vector $b$ by intersection
numbers on the special fibre:
\[
 A_{ij}=(Y_i\cdot Y_j\cdot H^{d-1}),\qquad
 b_i=(M\cdot Y_i\cdot H^{d-1}).                          \tag{E4.T.3}
\]
For $i\ne j$ one has $A_{ij}\geq0$, with strict inequality exactly when
$Y_i$ and $Y_j$ meet. This follows because their proper intersection is
an effective cycle and $H$ is ample on the projective special fibre.
The relation $\sum_jm_jY_j=\operatorname{div}(\varpi')$ gives $Am=0$.
The intersection graph is connected. Consequently
\[
 v^tAv=-\sum_{i<j}A_{ij}m_im_j
       \left(\frac{v_i}{m_i}-\frac{v_j}{m_j}\right)^2,
 \qquad \ker A=\mathbf Qm.                              \tag{E4.T.4}
\]
Constancy of intersection numbers in a flat projective family gives
\[
 \sum_i m_i b_i=(M_{K'}\cdot H_{K'}^{d-1})=0.
\]
Thus there is $a\in\mathbf Q^r$ with $Aa=-b$. Set
\[
 D=\sum_i a_iY_i,\qquad N=M+D\quad\text{in }\operatorname{Pic}(Y)_\mathbf Q.
                                                               \tag{E4.T.5}
\]
By construction $(N|_{Y_i}\cdot H|_{Y_i}^{d-1})=0$ for every $i$.
For $d=1$ this already says that $N|_{Y_i}$ is numerically trivial.
For $d\geq2$, the Hodge index theorem gives
\[
 (N|_{Y_i}^{\,2}\cdot H|_{Y_i}^{d-2})\leq0.             \tag{E4.T.6}
\]
Equality implies numerical triviality. The higher-dimensional version
follows by cutting with sufficiently ample integral hyperplanes and the
surface Hodge theorem; the equality implication is
\cite[Theorem~6.3(h), proof pp.~53--54, and Theorem~B.27, p.~79]{Kleiman2005}.
These statements apply to integral projective schemes, with no smoothness
hypothesis. Over a nonclosed field, pass to an algebraic closure: the
geometric components of an integral scheme are conjugate, and the
intersection numbers of descended bundles are equal on these components
up to their common positive multiplicity. The degree-zero condition,
the inequality, and its equality conclusion therefore give the same
statement over the original field. Denominators in $N$ are cleared before
applying the line-bundle theorem.

The same flat-family intersection calculation gives
\[
 \sum_i m_i(N|_{Y_i}^{\,2}\cdot H|_{Y_i}^{d-2})
       =(M_{K'}^{\,2}\cdot H_{K'}^{d-2})=0.              \tag{E4.T.7}
\]
Every summand is nonpositive and every $m_i$ is positive. Each is therefore
zero, and $N$ is numerically trivial on every special component, hence
on the whole special fibre. It is also numerically trivial on the generic
fibre. This construction gives a rational vertical correction; it does
not assert that an arbitrary generic line bundle extends integrally.

\paragraph{3. Galois symmetry removes the correction.}
Replace $D$ by $|G|^{-1}\sum_{g\in G}g^*D$. This is a rational average,
so no invertibility assumption on $|G|$ in $R$ is used. Since $M$ is pulled
back from $X$, it is $G$-invariant, and numerical triviality is preserved
by each (possibly semilinear over $R'$) automorphism. Thus $M+D$ remains
numerically trivial on both fibres.

Let $Z$ be the normalization of $X$ in $k(Y)$. It is finite over $X$ by
excellence. Normality of $Y$ factors $\pi$ as
\[
 Y\xrightarrow{h}Z\longrightarrow X,
\]
where $h$ is proper birational between normal schemes. The finite group
$G$ acts transitively on the prime divisors of $Z$ above $\xi$.
Here is the precise valuation argument. The DVR $\mathcal O_{X,\xi}$ has
a unique extension to the purely inseparable field $k(Y)^G$; the field
extension $k(Y)/k(Y)^G$ is Galois with group the faithful quotient of $G$,
and its Galois group permutes transitively the extensions of that valuation.
The same transitivity holds for their strict transforms on $Y$, since
a proper birational map to a normal scheme is an isomorphism at the
generic point of every target prime divisor.

These are exactly the components $Y_i$ that dominate $X_k$. Indeed a
vertical prime of the finite scheme $Z\to X$ has dimension $d$, hence
maps onto the integral $d$-dimensional special fibre. Their multiplicities
$m_i$ in (E4.T.2) are constant on the orbit: $G$ fixes the uniformizer of
$R$, whose divisor is the ramification index of $R'/R$ times (E4.T.2).
The averaged coefficients $a_i$ are constant on this orbit as well.
Consequently $a_i/m_i$ has a single value $c$ on every nonexceptional
vertical component. Subtract $c\operatorname{div}(\varpi')$ from $D$.
This is a principal rational divisor, so it preserves all numerical
conclusions and makes $D$ supported on $h$-exceptional prime divisors.

For any curve contracted by $h$, the degrees of $M$ and of $N=M+D$ are
zero. This includes curves over generic-fibre points, since there $D=0$
and $N=M$ is numerically trivial, as well as the special-fibre curves
treated above. Hence both $D$ and $-D$ are $h$-nef. The negativity lemma for proper
birational morphisms of normal schemes
\cite[Lemma~11.60, p.~436]{KollarFamilies} applied to both signs makes
$D\geq0$ and $D\leq0$, because $h_*D=0$. Thus $D=0$.
We conclude that $M|_{Y_{k'}}=\pi^*L|_{Y_{k'}}$ is numerically trivial.
Numerical triviality is detected by proper surjective pullback: for an
integral curve $C\subset X_k$, take a closed point of the generic fibre
of $Y_{k'}\times_{X_k}C\to C$ and close it to an integral curve $C'$.
The map $C'\to C$ is generically finite of positive degree, and the
projection formula gives
$0=\deg(M|_{C'})=[k(C'):k(C)]\deg(L|_C)$.
This proves the trait assertion.

\paragraph{4. From traits to an arbitrary base.}
First suppose $S$ is locally noetherian. The set (E4.T.1) is open by
the proper-flat Picard tau theorem
\cite[Theorem~2.1]{LaurentSchroer2024}; proper geometrically integral
fibres give universal $f_*\mathcal O_X=\mathcal O_S$.
Every specialization in a locally noetherian scheme is realized by a
map from the spectrum of a DVR (Stacks, Lemma~28.5.10, tag~054F).
The trait assertion shows that (E4.T.1) is stable under specialization.
Locally on the noetherian base it is constructible, being open; a
constructible specialization-stable subset is closed. Thus it is open
and closed.

For arbitrary $S$ the assertion is local, so take $S=\operatorname{Spec}A$.
Noetherian approximation descends $X/S$ and $L$ to some scheme
$S_i=\operatorname{Spec}A_i$ of finite type over $\mathbf Z$, with $X_i/S_i$
proper, flat and finitely presented. The needed exact limit results are
Stacks, Lemmas~32.10.1, 32.13.1, and 32.8.7 (tag~04AI), together with
finite-presentation descent of an invertible module. To ensure geometric
integrality of every fibre, let $E_i\subset S_i$ be the locus where the
fibre is geometrically reduced and geometrically irreducible. It is
constructible by Lemmas~37.26.5 and 37.27.7 (the latter tag~055B), and
$S\to S_i$ has image in $E_i$. Lemma~32.4.10, tag~05F4, shows that for
a sufficiently large index the entire $S_j\to S_i$ has image in $E_i$.
Replacing the model by $S_j$ therefore gives proper flat geometrically
integral fibres, while retaining $L_j$.
The cited limit and constructibility sources are collected in
\cite{BCEStacksTauDescent}.

The already proved noetherian assertion says $S_j^\tau(L_j)$ is open and
closed. Numerical triviality is invariant under field extension, so
$S^\tau(L)$ is exactly its inverse image. This proves (E4.T.1) for every
base. Finally, an fppf-local representative of a Picard section is an
actual invertible sheaf; applying (E4.T.1) on such a cover and descending
the open-and-closed locus proves the asserted closed immersion of the
tau subfunctor into the full Picard algebraic space. No schematicity
of the full Picard functor over $S$ is assumed or deduced.

\subsection{Fibrewise smoothness can improve at specialization}

\paragraph{Theorem E4.S.}
There is a projective flat family of geometrically integral reduced surfaces
over $R=\overline{\mathbf F}_2[[t]]$ whose geometric generic Picard
tau-part is nonsmooth, whereas its special Picard tau-part is smooth.
Thus the locus of fibrewise smoothness is not open. This is a precise
negative answer to the smoothness interpretation of 57:309(b$'$);
it does not identify fibrewise smoothness with relative smoothness.

Let $C\subset\mathbb P^2_R$ be the cubic
\[
 y^2z+txyz+tyz^2=x^3,\qquad o=[0:1:0].                 \tag{E4.S.1}
\]
Its discriminant and $c_4$ are
$\Delta=t^4(1+t^2)$ and $c_4=t^4$. Hence its generic fibre is a
smooth ordinary elliptic curve, with $j=t^8/(1+t^2)\ne0$.
The special fibre is the integral cuspidal cubic $y^2z=x^3$.
The coefficient of $y^2z$ makes $o$ smooth over $R$.
The involution
\[
 \sigma[x:y:z]=[x:y+tx+tz:z]                           \tag{E4.S.2}
\]
is inversion generically and fixes $o$.
Let $B/R$ be the constant ordinary elliptic curve
$y^2z+xyz=x^3+z^3$, let $P=(0,1)$ on its affine chart,
and let $\Gamma=(\mathbb Z/2\mathbb Z)_R$ act on $C\times_R B$
through $(\sigma,t_P)$. The nonzero point $P$ has order two, so this
action is free. Set
\[
 X=(C\times_R B)/\Gamma,\qquad Y=B/\langle P\rangle.
\]
The finite-group quotient is projective: an orbit tensor product of
an ample line bundle supplies an ample linearized line bundle.
Its quotient map is finite \'etale, so flatness descends, and all
geometric fibres of $X/R$ are integral and reduced. The invariant
section $o\times B$ gives a section $i:Y\to X$.

Write $G=\operatorname{Pic}^0_{C/R}$. It is a smooth separated group
scheme of finite presentation: the projective integral curve family
has its Picard scheme, and $H^2(C_s,\mathcal O_{C_s})=0$ removes
all infinitesimal obstructions. The degree-zero locus is its neutral
part and also its tau-part. The equality $\sigma^*=[-1]$ on $G$
holds generically; it holds everywhere because $G$ is flat over
the trait and the target is separated.
There is a base-change compatible isomorphism
\[
 \operatorname{Pic}^\tau_{X/R}
       \simeq Y^\vee\times_R G[2].                  \tag{E4.S.3}
\]
Here is the descent argument, including the singular special fibre.
Rigidification along $o$ makes a representing line bundle on $C$
unique up to a unique rigidified isomorphism. For a point of $G[2]$
the unique isomorphism with its $\sigma$-pullback respects this
rigidification, and its square is the identity. It therefore supplies
a canonical $\Gamma$-linearization. Pulling to $C\times B$ and
descending gives a line bundle on $X$ trivial on $i(Y)$.
Together with pullback from $Y$, this defines (E4.S.3).
Conversely, subtract the restriction to $i(Y)$ from a tau-trivial
line bundle on $X$. Its pullback to $C\times B$ is a bundle from
$C$ alone, by the product formula for tau and its trivialization on
$o\times B$. Descent makes its class $\sigma$-invariant. Its
linearization agrees with the canonical one: their ratio is a unit
on the proper connected product and is one along $o\times B$.
These constructions are inverse and commute with arbitrary base
change. Numerical triviality is detected by the finite surjective
cover, by the degree formula on curves. This proves (E4.S.3) as a
group-scheme assertion.

Over a geometric generic point,
$G[2]\simeq\mu_2\times(\mathbb Z/2\mathbb Z)$, so (E4.S.3)
is nonsmooth and has dimension one. On the special fibre,
$G_s\simeq\mathbb G_a$. Indeed, the cusp has normalization
$x=u^2,y=u^3$, conductor $(u^2)$, and the units quotient
\[
 (k[u]/(u^2))^\times/k^\times\simeq\mathbb G_a.
\]
The normalization has no degree-zero Picard group, giving the
assertion. Since $[2]$ is zero on $\mathbb G_a$,
\[
 (\operatorname{Pic}^\tau_{X/R})_s
       \simeq Y_s^\vee\times\mathbb G_a
\]
is smooth of dimension two. Its tangent dimension is two, as is
the generic tangent dimension, whereas the defect
$\dim\operatorname{Lie}-\dim$ drops from one to zero.
Consequently that defect is not upper semicontinuous either.
Relative smoothness remains an open condition; in this example it
fails at every special point because the relative dimension jumps.
The separate dimension function is upper semicontinuous, in agreement
with FGA VI Corollary~2.7, p.234.

The finite polynomial identities, Jacobian checks, discriminant,
and cusp parameterization are recorded in
\texttt{code/check\_cuspidal\_igusa.output.json}.
The singular-cubic extension of the rigidified descent proof received
a separate targeted AI audit from /root/dgj, which found no gap.

\subsection*{E4.Curve: schematicity for every smooth proper curve family}

\paragraph{Theorem E4.Curve.}
Let $C\to S$ be a proper smooth finitely presented morphism of schemes
with geometrically connected one-dimensional fibres, where $S$ is an
arbitrary scheme. Then the full fppf Picard functor $\operatorname{Pic}_{C/S}$
is represented by a smooth separated $S$-scheme. This settles the
geometrically normal curve case of the last question at 57:315. It does
not settle higher-dimensional families, or geometrically integral
singular curves without further hypotheses.

\paragraph{Proof.}
The arithmetic genus is locally constant, so we work on a locus where
it is a fixed integer $g$. If $g=0$, the dual of the relative canonical
bundle is relatively ample. If $g\geq2$, the relative canonical bundle
is relatively ample. In either case the family is locally projective
on $S$, and the projective Picard theorem gives a smooth separated
Picard scheme. These are exactly the arguments of
\cite[\S9.3, Theorem~1 and Remark~2, p.~252]{BCEBLR1990}.

Suppose $g=1$. Put $E=\operatorname{Pic}^0_{C/S}$; this is an abelian
scheme by \cite[\S9.4, Proposition~4, p.~260]{BCEBLR1990}.
The degree map decomposes the Picard algebraic space as
\[
 \operatorname{Pic}_{C/S}=\coprod_{n\in\mathbb Z}\operatorname{Pic}^n_{C/S}.
\]
The diagonal divisor defines $C\to\operatorname{Pic}^1_{C/S}$, and
this map is an isomorphism: after an etale cover giving a section of
$C/S$, it is the usual genus-one Abel map. Thus $C$ is an $E$-torsor.
For $n\neq0$, the map
\[
 C\longrightarrow\operatorname{Pic}^n_{C/S},\qquad
        x\longmapsto[\mathcal O_C(nx)]                  \tag{E4.Curve.1}
\]
is equivariant for multiplication by $n$ on $E$. After an etale
cover trivializing $C$, it is a translate of $[n]:E\to E$.
Consequently it is a finite locally free $E[n]$-torsor, and
\[
        \operatorname{Pic}^n_{C/S}=C/E[n]               \tag{E4.Curve.2}
\]
as fppf sheaves. This conclusion also holds when $n$ is divisible by
a residue characteristic; etaleness of $E[n]$ is unnecessary.

We show that (E4.Curve.2) is a scheme. First suppose $S$ is locally
noetherian. The assertion is Zariski local on $S$. Fix $s\in S$ and
choose an affine open $V\subset C$ meeting $C_s$. The image of $V$
in $S$ is open because $C/S$ is smooth. Replace $S$ by an affine open
neighborhood of $s$ in that image, and replace $V$ by its inverse
image. The latter is a quasi-affine scheme. Since $C$ is an $E$-torsor
and $V$ meets every fibre, the image of $E\times_S V\to C$ is all of
$C$. Apply \cite[\S6.6, Theorem~2(c), pp.~164--166]{BCEBLR1990}:
every finite set of points of $C$ lies in one affine open. In
particular, each finite $E[n]$-orbit does. The finite locally free
quotient theorem, \cite[Proposition~66.14.1, tag~07S6]{BCEStacksFiniteQuotient},
therefore represents $C/E[n]$ by a scheme. The degree-zero piece is
$E$ itself; hence every piece and their disjoint union are schemes.
Smoothness and separatedness follow etale locally from the projective
curve theorem, or from the torsor description.

Finally, the arbitrary-base assertion follows by approximation,
without assuming that schematicity descends along arbitrary fpqc
covers. Locally take $S$ affine. The proper smooth finitely presented
scheme $C/S$ descends to a proper smooth finitely presented scheme
$C_0/S_0$ with $S_0$ noetherian. Geometric connectedness of proper
smooth fibres is an open condition and the arithmetic genus is
locally constant. Enlarging the model and restricting to the
corresponding open locus therefore arranges geometrically connected
genus-one fibres. The noetherian argument makes
$\operatorname{Pic}_{C_0/S_0}$ a scheme; its base change to $S$ represents
$\operatorname{Pic}_{C/S}$ by the definition of the relative fppf
Picard functor. This completes the proof.

\paragraph{Provenance.}
The smooth-curve argument above was already given by K\k{e}stutis
\v{C}esnavi\v{c}ius in his answer of 28 April 2015 to Poonen's question
\cite{BCECesnavicius2015}. We independently checked the cited BLR and
finite-quotient statements before locating that answer. No novelty is
claimed for Theorem E4.Curve.

\paragraph{Why the genus-one torsor counterexample route fails.}
An arbitrary torsor under an abelian scheme can be a nonschematic
algebraic space. Such a torsor alone does not give the requested
counterexample, whose original family must be a scheme. If a genus-one
scheme family $C$ is used, its multiples in torsor cohomology are
exactly (E4.Curve.2), and the affine-neighborhood theorem above makes
all of them schemes. Thus taking a scheme genus-one torsor and hoping
that $\operatorname{Pic}^n$ becomes a nonschematic multiple cannot work.

\subsection{The two marginal notes on 57:310: mathematics and reading}

The multiplication margin is diplomatically transcribed as
``La multiplication par un multiple de $p$ est [surjective?]
[two illegible words] [ouverte?] (?) [comp. irr\'eductibles?]''.
The original image confirms ``Non, exemple'' and the branch diagram,
but does not justify removing the marked uncertainties. The mathematical
universal-openness question is nevertheless answered negatively by
Theorem E4.H, including for a smooth projective family. Its status is
therefore separate from the unresolved transcription of the qualifier.

The bracketed finite-type criterion has the alternatives ``au dessus''
and ``en dehors'' of an open set, with an uncertain dense/constructible
qualification. In the proper-flat-finitely-presented,
cohomologically-flat-in-degree-zero Picard setting over a noetherian
base, finite presentation of $\operatorname{Pic}^{\tau}$ already proves
that every such base change, open restriction, and closed complement has
finitely many irreducible components. Thus either indicated
finite-component condition holds in this Picard setting. This does not
prove a proposed general criterion for arbitrary group schemes.

\paragraph{An explicit warning about the abstract generalization.}
Over $S=\mathbf A^1_{\mathbf Q}$, enumerate the distinct points
$s_i=(t-i)$, $i\geq1$. Glue a copy $S_0$ of $S$ and copies $S_i$,
$i\geq1$, by identifying $S_i\setminus\{s_i\}$ with the corresponding
open of $S_0$. Call the resulting scheme $G$. Its fibre over $s_i$ has
two points, and every other fibre has one. It is an \'etale commutative
group scheme with the two-point fibre given the group law
$\mathbf Z/2$: on two different charts $S_i,S_j$, their sum is the
section that uses the extra origin at $s_i$ and $s_j$, obtained by gluing
on $S\setminus\{s_j\}$ and $S\setminus\{s_i\}$. On the same chart
the sum is the zero section. These formulas agree on chart intersections;
the group laws hold fibrewise and hence hold as morphisms into the
\'etale space.

The scheme $G$ is irreducible, since all its charts have the same dense
generic point, but is not quasi-compact: any finite collection of these
charts misses infinitely many extra origins. For every closed irreducible
$S'\subset S$, its restriction over a nonempty open of $S'$ has finitely
many irreducible components. Consequently this generic-open
component-count condition does not characterize finite type for arbitrary
locally finite-type group schemes. This is deliberately not presented as
a Picard realization. It only rules out silently replacing the uncertain
margin by that unrestricted abstract assertion.


\subsection*{E4.Pinching: an excluded route to nonschematic Picard spaces}

\paragraph{Proposition.}
Suppose $X'\to X$ is a finite schematically dominant surjection
of proper $S$-schemes, with its conductor square, and both conductor schemes are finite over
$S$. If $X'\to S$ is locally projective, so is $X\to S$.
Consequently, when $X/S$ is also flat, finitely presented and has
geometrically integral fibres, its full Picard functor is a scheme.

\paragraph{Proof.}
Work near any point $s\in S$ and choose a relatively ample invertible
sheaf $L'$ on $X'$. Its restriction to the inverse-image conductor is
trivial after shrinking $S$ around $s$: over the local scheme
$\operatorname{Spec}\mathcal O_{S,s}$, that conductor is finite and
hence semilocal, so every invertible module is free. The trivialization
spreads to a neighborhood by finite presentation. The conductor
fibre-product square therefore glues $L'$ to the trivial bundle on
the conductor of $X$, producing an invertible sheaf $L$ on $X$.
Since its pullback under the finite surjection is ample, $L$ is
relatively ample. This is the ampleness argument in
\cite[Proposition~2.4(i), p.~6]{BCEBrion2015}; for this argument only
finiteness of the conductors is needed. The final assertion is the
projective Picard theorem.

Thus finite-conductor pinching of a locally projective simultaneous
normalization, in particular of $\mathbf P^1_S$, cannot realize a
nonschematic Picard space. This does not assert that an arbitrary
pinching is flat or that normalization commutes with base change.
The general integral-curve schematicity question remains when no
such normalization over the original base is available. For normal
families of higher dimension, the arbitrary-base question likewise
remains unresolved here.

